\chapter{Mathematical foundations}
\label{ch:foundations}

This chapter sets up everything needed to read the main theorem of the book precisely: what a configuration is,
what ``density one'' and ``energy'' mean, which potentials are allowed, why Gaussians are the key case, and how a
single auxiliary function gives a lower bound on the energy of \emph{every} competitor. We give a complete human proof
of that bound for periodic configurations, explain how the bound passes from Gaussians to all completely monotone
potentials, and introduce just enough Lean~4 to read the formal statements. The chapter ends by reviewing a student
draft, the way a referee would.

\begin{warnbox}[: whose theorem is this?]
The main theorem (Theorem~\ref{thm:main-statement} below) is \emph{upstream's}: it was produced by an AI system and
released by OpenAI~\cite{oai} on 2026-09-23. It is not peer reviewed. Its Lean proof is kernel-checked and its
statement was accepted by the Comparator tool against the challenge statement in our pinned build. That is
\tierA{} for the \emph{formal statement}, conditional on the fidelity of the Lean definitions to the intended
mathematics. No human has audited that fidelity. Our own contribution, presented in Chapter~2, consists of corollaries
and an audit.
\end{warnbox}

Throughout, the tiers of the preface apply. In this chapter a \tierC{} label on a result that comes with a human
proof means that a human-readable argument is given here and has not been machine-checked. It does not mean the result
is doubtful.

% ====================================================================================================
\section{Point configurations, density and energy}
\label{sec:configs}

We work in the Euclidean plane $\R^2$ with the usual norm $|x|$. A \emph{configuration} is a subset $C\subseteq\R^2$.
Its points are the positions of identical particles. Because $C$ is a set, two particles never sit at the same place.

\begin{definition}[Local finiteness]\label{def:lf}
A configuration $C$ is \emph{locally finite} if $C\cap \overline{B}(0,R)$ is finite for every $R\in\R$, where
$\overline{B}(0,R)=\{x:|x|\le R\}$ is the closed disk of radius $R$ about the origin.
\end{definition}

For a locally finite $C$ we write $C_R = C\cap\overline B(0,R)$ for the finite set of points in the closed disk and
$n_C(R)=\#C_R$ for their number.

\begin{definition}[Centred density]\label{def:density}
A locally finite configuration $C$ has \emph{density one} if
\[
  \lim_{R\to\infty}\frac{n_C(R)}{\pi R^2}=1 .
\]
More generally it has density $\rho$ if the limit equals $\rho$.
\end{definition}

The density is \emph{centred}: it is measured in disks about the origin only, and the disks are \emph{closed}. Both
choices matter in principle. Centred density says nothing \emph{uniform}: a configuration of centred density one can
contain, far from the origin, unit disks holding $k^2$ points each (one such cluster near radius $2^k$ for each $k$),
because these add only $O((\log R)^3)$ points to the disk of radius $R$ \tierC{}. (For a \emph{fixed} other centre $c$,
the disks $\overline B(c,R)$ give the same limit, since
$\overline B(0,R-|c|)\subseteq\overline B(c,R)\subseteq\overline B(0,R+|c|)$.) And a point lying exactly on the circle $|x|=R$ is counted in the closed disk but not in the open
one. For locally finite configurations, the open-disk and closed-disk density and lower energy coincide
\tierC{} (human proof in the preregistration of D4). For the formal status of this fact: \DFourStatus.

\begin{leanbox}[: upstream's definitions of local finiteness, disk points, count and density]
Quoted verbatim from \lean{lean/ComparatorChallenges/TriangularEnergy.lean} (openai/math \lean{adc7f124}); the
empty-set symbol is typeset through an escape because the book's listing setup has no glyph for it.
\begin{lstlisting}[language=Lean,escapeinside={(*}{*)}]
abbrev Plane := EuclideanSpace ℝ (Fin 2)

def LocallyFinite (C : Set Plane) : Prop :=
  ∀ R : ℝ, (C ∩ Metric.closedBall (0 : Plane) R).Finite

def diskPoints (C : Set Plane) (R : ℝ) : Finset Plane :=
  if h : (C ∩ Metric.closedBall (0 : Plane) R).Finite then h.toFinset else (*$\emptyset$*)

def diskCount (C : Set Plane) (R : ℝ) : ℕ := (diskPoints C R).card

def DensityOne (C : Set Plane) : Prop :=
  Tendsto (fun R : ℝ => (diskCount C R : ℝ) / (Real.pi * R ^ 2))
    atTop (𝓝 1)
\end{lstlisting}
\end{leanbox}

Three details of the Lean text deserve attention, because a reader of the human definition could miss them.
\begin{itemize}
\item \lean{diskPoints} returns the empty set when $C_R$ is infinite. The definition is therefore total, and every
theorem that needs the count to mean something must assume \lean{LocallyFinite C}.
\item \lean{Tendsto ... atTop} $(\mathcal N\,1)$ is Lean's way of writing $\lim_{R\to\infty}=1$, with $R$ a \emph{real} parameter.
It is not a limit along integers.
\item At $R=0$ the denominator $\pi R^2$ is zero, and Lean's convention $x/0=0$ makes the quotient $0$. This is
harmless because only large $R$ matter.
\end{itemize}

\subsection{Pair potentials and the lower energy}

A \emph{pair potential} is written as a function $g$ of the \emph{squared} distance $t=|x-y|^2$. This convention, used by
CKMRV~\cite{ckmrv} and by upstream~\cite{oai}, makes Gaussians $e^{-\alpha t}$ the basic building blocks. The Riesz
potential $|x-y|^{-s}$ becomes $g(t)=t^{-s/2}$.

\begin{definition}[Disk energy and lower energy]\label{def:energy}
Let $g:(0,\infty)\to[0,\infty)$ and let $C$ be locally finite. The \emph{disk energy} at radius $R$ is
\[
  E_g(C;R)=\frac{1}{n_C(R)}\sum_{x\in C_R}\ \sum_{\substack{y\in C_R\\ y\ne x}} g\bigl(|x-y|^2\bigr)\ \in[0,\infty],
\]
with the convention $1/0=\infty$ in $[0,\infty]$ (and $\infty\cdot 0=0$). The \emph{lower energy} of $C$ is
\[
  E_g(C)=\liminf_{R\to\infty}E_g(C;R)\ \in[0,\infty].
\]
\end{definition}

Each unordered pair $\{x,y\}$ is counted twice, as $(x,y)$ and $(y,x)$. This is the ``ordered-pair'' normalisation:
for a lattice it gives $\sum_{a\ne0}g(|a|^2)$, the energy seen by one particle, without a factor $\tfrac12$.

\begin{leanbox}[: upstream's disk energy and energy]
Quoted verbatim from \lean{TriangularEnergy.lean}.
\begin{lstlisting}[language=Lean]
def diskEnergy (g : ℝ → ℝ) (C : Set Plane) (R : ℝ) : ℝ≥0∞ :=
  (diskCount C R : ℝ≥0∞)⁻¹ *
    ∑ x ∈ diskPoints C R, ∑ y ∈ (diskPoints C R).erase x,
      ENNReal.ofReal (g (‖x - y‖ ^ 2))

def energy (g : ℝ → ℝ) (C : Set Plane) : ℝ≥0∞ :=
  Filter.liminf (diskEnergy g C) atTop
\end{lstlisting}
\end{leanbox}

The values live in $\mathbb R_{\ge0}\cup\{\infty\}$ (Lean: \lean{ENNReal}), the extended non-negative reals $[0,\infty]$. Then every energy is defined, possibly
infinite, and no summability hypothesis is needed to \emph{state} the theorem. The function \lean{ENNReal.ofReal} sends
negative reals to $0$, which is irrelevant for the non-negative potentials used here.

\begin{warnbox}[: lower energy, not energy]
$E_g(C)$ is a $\liminf$. CKMRV's Definition~1.3~\cite{ckmrv} of universal optimality also requires the limit to
\emph{exist} for the optimal configuration. The Lean notion used in this book uses only the lower energy, and is
therefore weaker than CKMRV's \tierC{}. When we say ``energy'' below we always mean the lower energy unless the limit is
shown to exist.
\end{warnbox}

\begin{example}[Why the liminf]\label{ex:liminf}
The $\liminf$ is the smallest value that the disk energies approach infinitely often. A competitor whose disk energies
oscillate, for example because of rings of extra points at sparse radii, is judged by its \emph{best} radii. A lower
bound $\liminf\ge L$ is therefore a strong kind of lower bound: for every $\varepsilon>0$, all sufficiently large radii
have disk energy at least $L-\varepsilon$. \tierC{}
\end{example}

% ====================================================================================================
\section{Lattices and the triangular lattice}
\label{sec:lattices}

A \emph{lattice} in $\R^2$ is a set $\Lambda=\{jv_1+kv_2:(j,k)\in\Z^2\}$ for linearly independent $v_1,v_2$. Its
\emph{covolume} $|\Lambda|=|\det(v_1,v_2)|$ is the area of a fundamental parallelogram. A lattice has density
$1/|\Lambda|$, so a density-one lattice has covolume one.

\begin{definition}[The triangular lattice $\Lat$]\label{def:A}
Let $b=\sqrt3/2$ and
\[
  p(j,k)=\frac{1}{\sqrt b}\Bigl(j+\frac k2,\ k\,b\Bigr),\qquad (j,k)\in\Z^2,
  \qquad \Lat=\{p(j,k):(j,k)\in\Z^2\}.
\]
\end{definition}

\begin{leanbox}[: upstream's lattice]
Quoted verbatim from \lean{TriangularEnergy.lean}.
\begin{lstlisting}[language=Lean]
def b : ℝ := Real.sqrt 3 / 2

def triangularPoint (jk : ℤ × ℤ) : Plane :=
  (Real.sqrt b)⁻¹ •
    (EuclideanSpace.single 0 ((jk.1 : ℝ) + (jk.2 : ℝ) / 2) +
      EuclideanSpace.single 1 ((jk.2 : ℝ) * b))

def A : Set Plane := Set.range triangularPoint

/-- The full nonzero-lattice energy, allowing a divergent series. -/
def latticeEnergy (g : ℝ → ℝ) : ℝ≥0∞ :=
  ∑' a : {x : Plane // x ∈ A ∧ x ≠ 0}, ENNReal.ofReal (g (‖a.val‖ ^ 2))
\end{lstlisting}
\end{leanbox}

The basis vectors are $v_1=p(1,0)=b^{-1/2}(1,0)$ and $v_2=p(0,1)=b^{-1/2}(\tfrac12,b)$.

\begin{proposition}[Covolume and Gram form]\label{prop:gram}
$|\det(v_1,v_2)|=1$, so $\Lat$ has covolume one. The Gram matrix is
\[
  G=\begin{pmatrix} v_1\cdot v_1 & v_1\cdot v_2\\ v_2\cdot v_1 & v_2\cdot v_2\end{pmatrix}
   =\frac1b\begin{pmatrix}1&\tfrac12\\ \tfrac12&1\end{pmatrix},
\]
and consequently
\begin{equation}\label{eq:normsq}
  |p(j,k)|^2=\frac1b\,(j^2+jk+k^2)=\frac{2}{\sqrt3}\,(j^2+jk+k^2).
\end{equation}
\end{proposition}
\begin{humanproof}
$\det(v_1,v_2)=b^{-1}\det\begin{pmatrix}1&\tfrac12\\0&b\end{pmatrix}=b^{-1}\cdot b=1$. Next,
$v_1\cdot v_1=b^{-1}$, $v_1\cdot v_2=b^{-1}\cdot\tfrac12$ and $v_2\cdot v_2=b^{-1}(\tfrac14+b^2)=b^{-1}(\tfrac14+\tfrac34)=b^{-1}$.
Finally $|jv_1+kv_2|^2=(j,k)\,G\,(j,k)^{\mathsf T}=b^{-1}(j^2+jk+k^2)$. \tierC{} for the covolume and the Gram
matrix (human proof). The norm identity~\eqref{eq:normsq} is also kernel-checked \tierA{}.
\end{humanproof}

\begin{leanbox}[: the norm identity (our file \lean{formal_cert/Triangular/RieszFinite.lean})]
\begin{lstlisting}[language=Lean]
/-- The squared norm of a lattice point. -/
lemma norm_sq_triangularPoint (j k : ℤ) :
    ‖triangularPoint (j, k)‖ ^ 2 = (1 / b) * ((j : ℝ) ^ 2 + (j : ℝ) * k + (k : ℝ) ^ 2) := by
\end{lstlisting}
\end{leanbox}

All six nearest neighbours of the origin are at squared distance $1/b=2/\sqrt3$. The minimal distance of the
density-one triangular lattice is therefore $(2/\sqrt3)^{1/2}\approx1.0746$ \tierC{}. In the physics chapters the same
lattice appears with spacing $a=\sqrt{2/\sqrt3}$.

\subsection{The dual lattice}

\begin{definition}[Dual lattice]
The dual of a lattice $\Lambda$ is $\Lambda^*=\{w\in\R^2: w\cdot\lambda\in\Z\text{ for all }\lambda\in\Lambda\}$.
Its Gram matrix in the dual basis is $G^{-1}$, and $|\Lambda^*|=1/|\Lambda|$.
\end{definition}

\begin{proposition}[$\Lat^*$ is a rotated copy of $\Lat$]\label{prop:dual}
Let $J$ be the rotation by $90^\circ$. Then $\Lat^*=J\Lat$. In particular $\Lat^*$ is isometric to $\Lat$ and has
covolume one.
\end{proposition}
\begin{humanproof}
$Jv_1=b^{-1/2}(0,1)$ and $Jv_2=b^{-1/2}(-b,\tfrac12)$. Compute $Jv_1\cdot v_1=0$, $Jv_1\cdot v_2=b^{-1}b=1$,
$Jv_2\cdot v_1=-b^{-1}b=-1$ and $Jv_2\cdot v_2=b^{-1}(-\tfrac b2+\tfrac b2)=0$. All four are integers, so
$J\Lat\subseteq\Lat^*$. Both lattices have covolume $1$: $J$ preserves area, and $|\Lat^*|=1/|\Lat|=1$. A sublattice of
the same covolume is the whole lattice, so $J\Lat=\Lat^*$. As a cross-check, $G^{-1}=b^{-1}\begin{pmatrix}1&-\frac12\\-\frac12&1\end{pmatrix}$
(note $\det G=b^{-2}\cdot\tfrac34=1$), whose quadratic form $b^{-1}(j^2-jk+k^2)$ becomes $b^{-1}(j^2+jk+k^2)$ under
$k\mapsto-k$. \tierC{}
\end{humanproof}

\subsection{Lattice energy as an Epstein-type sum}

For a lattice the natural energy is the energy seen from one lattice point:
\[
  \mathcal E_g(\Lat)=\sum_{a\in\Lat\setminus\{0\}}g\bigl(|a|^2\bigr)
  =\sum_{(j,k)\ne(0,0)} g\Bigl(\tfrac{2}{\sqrt3}(j^2+jk+k^2)\Bigr)\in[0,\infty],
\]
which is upstream's \lean{latticeEnergy}. Group the terms by the value $n=j^2+jk+k^2$, and let $r(n)$ be the number of
$(j,k)$ with that value. Then $\mathcal E_g(\Lat)=\sum_{n\ge1}r(n)\,g(2n/\sqrt3)$. For the Riesz potential
$g(t)=t^{-s/2}$ this is
\[
  \mathcal E_{t^{-s/2}}(\Lat)=b^{s/2}\sum_{(j,k)\ne0}(j^2+jk+k^2)^{-s/2},
\]
an Epstein zeta function of the form $j^2+jk+k^2$ evaluated at $s/2$. It is finite exactly when $s>2$. Finiteness for
$s>2$ is kernel-checked as \lean{latticeEnergy_riesz_lt_top} \tierA{}. Divergence for $s\le2$ follows by comparing the
sum with $\int_{|x|\ge1}|x|^{-s}\,dx$ \tierC{}.

\begin{example}[The first shells]\label{ex:shells}
The smallest values of $j^2+jk+k^2$ and their multiplicities can be checked by hand \tierC{}.
\begin{itemize}
\item $n=1$: $(\pm1,0)$, $(0,\pm1)$, $\pm(1,-1)$, so $r(1)=6$.
\item $n=3$: $\pm(1,1)$, $\pm(2,-1)$, $\pm(1,-2)$, so $r(3)=6$.
\item $n=4$: twice the $n=1$ vectors, so $r(4)=6$.
\item $n=7$: the twelve vectors $\pm(2,1)$, $\pm(1,2)$, $\pm(3,-1)$, $\pm(3,-2)$, $\pm(1,-3)$, $\pm(2,-3)$, so $r(7)=12$.
\item $n=2,5,6$ do not occur.
\end{itemize}
The script below cross-checks these counts \tierX{}. Running it printed the line shown.
\begin{lstlisting}[language=Python]
from collections import Counter

M = 20
count = Counter(j * j + j * k + k * k
                for j in range(-M, M + 1) for k in range(-M, M + 1))
print([(n, count[n]) for n in range(1, 14) if count[n] > 0])
\end{lstlisting}
\begin{lstlisting}
[(1, 6), (3, 6), (4, 6), (7, 12), (9, 6), (12, 6), (13, 12)]
\end{lstlisting}
For the Gaussian $g(t)=e^{-\pi t}$ (that is, $\alpha=1$ in the normalisation of \S\ref{sec:lp}), the shell
contributions $r(n)e^{-\pi n/b}$ decay very fast. The following script was run and printed the output shown \tierX{}.
\begin{lstlisting}[language=Python]
import math

b = math.sqrt(3) / 2
shells = [(1, 6), (3, 6), (4, 6), (7, 12)]
for n, m in shells:
    print(n, m, m * math.exp(-math.pi * n / b))
print(sum(m * math.exp(-math.pi * n / b) for n, m in shells))
\end{lstlisting}
\begin{lstlisting}
1 6 0.15947960085851695
3 6 0.00011267119989805458
4 6 2.9947929979986557e-06
7 12 1.1247564017693356e-10
0.15959526696388865
\end{lstlisting}
So $\sum_{a\in\Lat\setminus0}e^{-\pi|a|^2}\approx0.1595953$, and the first shell alone carries more than $99.9\%$ of it
\tierX{}. Compare this with the full truncated theta sum in Exercise~\ref{ex:py-theta}.
\end{example}

% ====================================================================================================
\section{Completely monotone potentials}
\label{sec:cm}

\begin{definition}[Completely monotone; admissible potential]\label{def:cm}
A function $g:(0,\infty)\to\R$ is \emph{completely monotone} if it is $C^\infty$ on $(0,\infty)$ and
\[
  (-1)^r g^{(r)}(t)\ \ge\ 0\qquad\text{for all } r=0,1,2,\dots\ \text{and all } t>0.
\]
Upstream calls $g$ \emph{admissible} if it is $C^\infty$ on $(0,\infty)$, non-negative there, and completely monotone.
The non-negativity is the case $r=0$.
\end{definition}

\begin{leanbox}[: upstream's admissible potentials]
Quoted verbatim from \lean{TriangularEnergy.lean}.
\begin{lstlisting}[language=Lean]
def AdmissiblePotential (g : ℝ → ℝ) : Prop :=
  ContDiffOn ℝ (⊤ : ℕ∞) g (Set.Ioi 0) ∧
  (∀ t : ℝ, 0 < t → 0 ≤ g t) ∧
  ∀ (r : ℕ) (t : ℝ), 0 < t → 0 ≤ (-1 : ℝ) ^ r * iteratedDeriv r g t
\end{lstlisting}
\end{leanbox}

Note three things.
\begin{enumerate}
\item The inequalities are \emph{weak}: $g\equiv1$ is admissible, and so is $g\equiv0$. A definition with strictly
alternating signs would be a different (smaller) class.
\item The smoothness is $C^\infty$ (\lean{ContDiffOn} with smoothness index $\top$). It is \emph{not} real-analyticity. Analyticity is a
consequence, through Bernstein's theorem below, but it is not part of the definition.
\item Only $t>0$ matters. The value of $g$ at $0$ (or at negative $t$) never enters the energy, because distinct points
have $|x-y|^2>0$.
\end{enumerate}

\begin{example}[Gaussians]\label{ex:gauss}
For $\alpha\ge0$, $g(t)=e^{-\alpha t}$ satisfies $(-1)^rg^{(r)}(t)=\alpha^re^{-\alpha t}\ge0$. So it is admissible
\tierC{}. In space, $g(|x|^2)=e^{-\alpha|x|^2}$ is a Gaussian. The case $\alpha=0$ is the constant potential.
\end{example}

\begin{example}[Riesz potentials]\label{ex:riesz}
For $s>0$, $g(t)=t^{-s/2}$ has $(-1)^rg^{(r)}(t)=\tfrac s2(\tfrac s2+1)\cdots(\tfrac s2+r-1)\,t^{-s/2-r}>0$. It is
admissible, and this is kernel-checked as \lean{TriangularRiesz.riesz_admissible} \tierA{}. The Lean proof rewrites
the $r$-th derivative with \lean{Real.iter_deriv_rpow_const}, converts the falling factorial at $-s/2$ into a rising
factorial at $s/2$ (\lean{ascPochhammer_eval_neg_eq_descPochhammer}), and uses positivity of the rising factorial
(\lean{ascPochhammer_pos}). At $s=0$ it is the constant $1$, which is admissible. For $s<0$ it is not: for instance
$t\mapsto t$ (the case $s=-2$) has first derivative $1>0$. That case is kernel-checked as
\lean{riesz_neg_two_not_admissible} (statement: not \lean{AdmissiblePotential (TriangularRiesz.riesz (-2))}) \tierA{}.
\end{example}

\begin{example}[Yukawa in the variable $t=r^2$]\label{ex:yukawa}
The screened Coulomb (Yukawa) potential $V(r)=e^{-\kappa r}/r$, $\kappa>0$, becomes
$g(t)=e^{-\kappa\sqrt t}/\sqrt t$. We claim it is completely monotone \tierC{}. The argument uses the classical integral
\begin{equation}\label{eq:yukawa-laplace}
  \frac{e^{-\kappa\sqrt t}}{\sqrt t}=\frac{1}{\sqrt\pi}\int_0^\infty e^{-tu}\,u^{-1/2}e^{-\kappa^2/(4u)}\,du\qquad(t>0).
\end{equation}
It follows from $\int_0^\infty u^{-1/2}e^{-tu-a/u}\,du=\sqrt{\pi/t}\,e^{-2\sqrt{at}}$ with $a=\kappa^2/4$.
Formula~\eqref{eq:yukawa-laplace} writes $g$ as a positive mixture of exponentials $e^{-tu}$. Differentiating under the
integral (dominated convergence, since $u^re^{-tu}$ is integrable against the weight for $t>0$) gives
$(-1)^rg^{(r)}(t)=\pi^{-1/2}\int_0^\infty u^r e^{-tu}u^{-1/2}e^{-\kappa^2/(4u)}\,du\ge0$. A numerical spot check of
\eqref{eq:yukawa-laplace} at $\kappa=2$ ran as follows \tierX{}:
\begin{lstlisting}[language=Python]
import math
from scipy.integrate import quad

kappa = 2.0
for t in (0.25, 1.0, 4.0):
    lhs = math.exp(-kappa * math.sqrt(t)) / math.sqrt(t)
    integrand = lambda u: u ** -0.5 * math.exp(-kappa ** 2 / (4 * u)) * math.exp(-t * u)
    val, err = quad(integrand, 0, math.inf)
    print(f"t={t:4.2f}  lhs={lhs:.12f}  rhs={val / math.sqrt(math.pi):.12f}")
\end{lstlisting}
\begin{lstlisting}
t=0.25  lhs=0.735758882343  rhs=0.735758882343
t=1.00  lhs=0.135335283237  rhs=0.135335283237
t=4.00  lhs=0.009157819444  rhs=0.009157819444
\end{lstlisting}
The spot check steers but does not prove. The proof is the integral identity. No Lean proof of Yukawa admissibility
exists in this book's files.
\end{example}

Two closure properties are useful \tierC{}. Positive combinations of completely monotone functions are completely
monotone, which is obvious from the definition. Products of completely monotone functions are completely monotone,
because Leibniz's rule gives $(-1)^r(fg)^{(r)}=\sum_i\binom ri\,[(-1)^if^{(i)}]\,[(-1)^{r-i}g^{(r-i)}]\ge0$.

\begin{theorem}[Bernstein's theorem]\label{thm:bernstein}
A function $g:(0,\infty)\to\R$ is completely monotone if and only if there is a non-negative Borel measure $\mu$ on
$[0,\infty)$ such that
\[
  g(t)=\int_{[0,\infty)}e^{-\alpha t}\,d\mu(\alpha)\qquad\text{for all }t>0,
\]
the integral being finite for every $t>0$. The measure $\mu$ is then unique. \tierL{} (classical result of Bernstein,
in the form attributed by upstream~\cite{oai} to Bernstein--Hausdorff--Widder. Not proved in this book.)
\end{theorem}

Two features of the statement matter for energy problems.
\begin{itemize}
\item \textbf{The measure lives on $[0,\infty)$, closed at $0$.} An atom $\mu(\{0\})=c>0$ contributes the constant $c$
to $g$. For instance $g(t)=1+t^{-2}$ is completely monotone with $\mu=\delta_0+(\text{a density on }(0,\infty))$.
\item \textbf{A constant potential has infinite energy at density one.} If $g\equiv1$ then $E_1(C;R)=n_C(R)-1$ (each of
the $n_C(R)$ points sees $n_C(R)-1$ others), and $n_C(R)\to\infty$ at density one. Hence $E_1(C)=\infty$ for every
density-one $C$ \tierC{} (the human proof is this one line). Upstream proves the same fact in Lean as
\lean{energy_one} (file \lean{OAI/Analysis/Triangular/Energy/Potentials.lean}), with the formal statement
\lean{energy (fun _ => 1) C} $=\infty$ under \lean{DensityOne C}. It is part of the build whose main theorem is kernel-checked.
The lattice energy of the constant is also $\infty$, so the atom at $0$ never breaks the inequality: both sides are
$\infty$.
\end{itemize}

% ====================================================================================================
\section{Fourier analysis and Poisson summation}
\label{sec:fourier}

We use the normalisation of upstream~\cite{oai}:
\[
  \widehat f(\xi)=\int_{\R^2}f(x)\,e^{-2\pi i\,x\cdot\xi}\,dx .
\]
With this choice $\widehat{\widehat f}(x)=f(-x)$ for Schwartz $f$ (Fourier inversion, \tierL{}, classical), and no
factors of $2\pi$ appear in the Poisson summation formula.

\begin{proposition}[Gaussians]\label{prop:gauss-ft}
For $\alpha>0$ let $G_\alpha(x)=e^{-\pi\alpha|x|^2}$ on $\R^2$. Then
\[
  \widehat{G_\alpha}(\xi)=\frac1\alpha\,e^{-\pi|\xi|^2/\alpha}.
\]
In particular $\widehat{G_\alpha}>0$, and $G_1=\widehat{G_1}$.
\end{proposition}
\begin{humanproof}
In one variable, $\int_\R e^{-\pi\alpha x^2}e^{-2\pi ix\xi}\,dx=\alpha^{-1/2}e^{-\pi\xi^2/\alpha}$. This is the standard
computation: complete the square and shift the contour, or check that both sides solve the same first-order ODE in
$\xi$ with the same value at $\xi=0$. In two variables the integral factorises into a product of two such integrals,
giving $\alpha^{-1/2}\cdot\alpha^{-1/2}=\alpha^{-1}$. \tierC{}
\end{humanproof}

\begin{theorem}[Poisson summation for a lattice]\label{thm:poisson}
Let $\Lambda\subset\R^2$ be a lattice of covolume $|\Lambda|$ and $f$ a Schwartz function on $\R^2$. For every
$v\in\R^2$,
\[
  \sum_{\lambda\in\Lambda}f(\lambda+v)=\frac{1}{|\Lambda|}\sum_{w\in\Lambda^*}\widehat f(w)\,e^{2\pi i\,w\cdot v},
\]
both sides converging absolutely. \tierL{} (classical. The sketch below is \tierC{}.)
\end{theorem}
\begin{proof}[Sketch for Schwartz functions]
Let $F(v)=\sum_{\lambda}f(\lambda+v)$. Since $f$ and all its derivatives decay faster than any power, the series and all
its term-by-term derivatives converge uniformly on compact sets. So $F$ is $C^\infty$, and it is $\Lambda$-periodic.
The characters of $\R^2/\Lambda$ are $e^{2\pi iw\cdot v}$ with $w\in\Lambda^*$. The Fourier coefficient of $F$ is
\[
  c_w=\frac1{|\Lambda|}\int_{D}F(v)e^{-2\pi iw\cdot v}\,dv=\frac1{|\Lambda|}\int_{\R^2}f(x)e^{-2\pi iw\cdot x}\,dx
  =\frac{\widehat f(w)}{|\Lambda|},
\]
where $D$ is a fundamental domain. We unfold the integral over $D$ and the sum over $\lambda$ into one integral over
$\R^2$, using $e^{-2\pi iw\cdot\lambda}=1$. A smooth periodic function has rapidly decaying Fourier coefficients and
equals its absolutely convergent Fourier series everywhere. This gives the formula.
\end{proof}

\begin{example}[Theta functions of $\Lat$]\label{ex:theta}
Apply Theorem~\ref{thm:poisson} to $f=G_\alpha$, $\Lambda=\Lat$ and $v=0$. Proposition~\ref{prop:gauss-ft} gives
$\widehat{G_\alpha}$. By Proposition~\ref{prop:dual}, $|w|$ runs over the same values on $\Lat^*$ as $|a|$ on $\Lat$, with
the same multiplicities. Writing $\theta_\Lat(\alpha)=\sum_{a\in\Lat}e^{-\pi\alpha|a|^2}$ we get the functional equation
\[
  \theta_\Lat(\alpha)=\frac1\alpha\,\theta_\Lat(1/\alpha)\qquad(\alpha>0).
\]
\tierC{} (it follows from the two propositions and the theorem). Exercise~\ref{ex:py-theta} checks it in floating
point.
\end{example}

% ====================================================================================================
\section{The linear-programming bound}
\label{sec:lp}

The idea goes back to Cohn and Elkies for sphere packing and to Cohn and Kumar~\cite{ck} for energy. It turns the
problem ``bound the energy of every configuration'' into the problem ``find one good auxiliary function''. We state it
for the Gaussian $g_\alpha(t)=e^{-\pi\alpha t}$, so that $g_\alpha(|x|^2)=G_\alpha(x)$.

\begin{definition}[Auxiliary function]\label{def:aux}
Let $\alpha>0$. A radial Schwartz function $f:\R^2\to\R$ is an \emph{auxiliary function for $G_\alpha$} if
\begin{enumerate}
\item[(i)] $f(x)\le G_\alpha(x)$ for all $x\ne0$ ($f$ is a \emph{minorant} of the potential), and
\item[(ii)] $\widehat f(\xi)\ge0$ for all $\xi$.
\end{enumerate}
It is \emph{sharp for $\Lat$} if in addition
\begin{enumerate}
\item[(iii)] $f(a)=G_\alpha(a)$ for all $a\in\Lat\setminus\{0\}$, and
\item[(iv)] $\widehat f(w)=0$ for all $w\in\Lat^*\setminus\{0\}$.
\end{enumerate}
\end{definition}

Since $f$ is radial and real, $\widehat f$ is real and radial, so (ii) makes sense. Conditions (iii) and (iv) are the
``root conditions''. In (iii), since $G_\alpha-f\ge0$, every point of $\Lat\setminus0$ is a minimum of $G_\alpha-f$.

\subsection{Periodic configurations: a complete proof}

\begin{definition}[Periodic configuration and its energy]\label{def:periodic}
Let $\Lambda$ be a lattice and $v_1,\dots,v_N\in\R^2$ with $v_i-v_j\notin\Lambda$ for $i\ne j$. The configuration
$C=\bigcup_{i=1}^N(\Lambda+v_i)$ is \emph{$\Lambda$-periodic with $N$ points per period}. Its density is
$N/|\Lambda|$. Its \emph{periodic energy} is
\[
  E^{\mathrm{per}}_g(C)=\frac1N\sum_{i=1}^N\sum_{j=1}^N\ \sum_{\substack{\lambda\in\Lambda\\ \lambda+v_j-v_i\ne0}}
  g\bigl(|\lambda+v_j-v_i|^2\bigr),
\]
the average over the $N$ cosets of the energy seen by one particle.
\end{definition}

The condition $\lambda+v_j-v_i\ne0$ excludes exactly one term: $i=j$ and $\lambda=0$. For $i\ne j$ the vector
$\lambda+v_j-v_i$ is never zero, because the cosets are distinct.

\begin{theorem}[LP bound, periodic case]\label{thm:lp-periodic}
Let $\alpha>0$, let $f$ be an auxiliary function for $G_\alpha$, and let $C$ be a $\Lambda$-periodic configuration of
density one ($|\Lambda|=N$). Then
\[
  E^{\mathrm{per}}_{g_\alpha}(C)\ \ge\ \widehat f(0)-f(0).
\]
If $f$ is sharp for $\Lat$, then $E^{\mathrm{per}}_{g_\alpha}(\Lat)=\widehat f(0)-f(0)=\sum_{a\in\Lat\setminus0}e^{-\pi\alpha|a|^2}$.
Hence $\Lat$ minimises the periodic Gaussian energy among periodic density-one configurations. \tierC{} (complete
human proof below.)
\end{theorem}
\begin{humanproof}
All series below converge absolutely, because $G_\alpha$ and $f$ are Schwartz.

\emph{Step 1 (minorant).} By (i), each term of $E^{\mathrm{per}}$ is at least the corresponding value of $f$:
\[
  N\,E^{\mathrm{per}}_{g_\alpha}(C)=\sum_{i,j}\sum_{\lambda+v_j-v_i\ne0}G_\alpha(\lambda+v_j-v_i)
  \ \ge\ \sum_{i,j}\sum_{\lambda+v_j-v_i\ne0}f(\lambda+v_j-v_i).
\]
\emph{Step 2 (put back the diagonal).} The excluded terms are exactly the $N$ terms with $i=j$ and $\lambda=0$, each
equal to $f(0)$:
\[
  \sum_{i,j}\sum_{\lambda+v_j-v_i\ne0}f(\lambda+v_j-v_i)=\sum_{i,j}\sum_{\lambda\in\Lambda}f(\lambda+v_j-v_i)\;-\;N f(0).
\]
\emph{Step 3 (Poisson).} By Theorem~\ref{thm:poisson} with $v=v_j-v_i$ and $|\Lambda|=N$,
\[
  \sum_{i,j}\sum_{\lambda}f(\lambda+v_j-v_i)=\frac1N\sum_{w\in\Lambda^*}\widehat f(w)\sum_{i,j}e^{2\pi i\,w\cdot(v_j-v_i)}
  =\frac1N\sum_{w\in\Lambda^*}\widehat f(w)\,\Bigl|\sum_{j=1}^Ne^{2\pi i\,w\cdot v_j}\Bigr|^2 .
\]
The interchange of the finite sum over $(i,j)$ with the sum over $w$ is harmless.\\
\emph{Step 4 (positivity).} Each term on the right is $\ge0$ by (ii). Keeping only $w=0$, where the inner sum equals
$N$, gives a lower bound of $\frac1N\widehat f(0)N^2=N\widehat f(0)$.

Combining the four steps, $N\,E^{\mathrm{per}}\ge N\widehat f(0)-Nf(0)$. Divide by $N$.

\emph{Equality on $\Lat$.} Take $\Lambda=\Lat$, $N=1$, $v_1=0$. In Step~1, (iii) makes every term an equality. In
Step~4, (iv) makes every dropped term zero. So both inequalities are equalities:
$\sum_{a\ne0}G_\alpha(a)=\widehat f(0)-f(0)$. Since $G_\alpha(a)=e^{-\pi\alpha|a|^2}$, this is the stated value.
\end{humanproof}

\begin{remark}[Other densities]
The same proof at density $\rho=N/|\Lambda|$ gives $E^{\mathrm{per}}\ge\rho\,\widehat f(0)-f(0)$. The upstream theorem
is stated at density one. Chapter~2 transfers it to every density by scaling. \tierC{}
\end{remark}

\subsection{From the periodic energy to the lower energy}

Theorem~\ref{thm:lp-periodic} uses $E^{\mathrm{per}}$, not the disk energy of Definition~\ref{def:energy}. These are
different definitions, and we connect them explicitly rather than silently identifying them.

\begin{proposition}[Disk energy of a periodic configuration]\label{prop:periodic-limit}
Let $C$ be periodic as in Definition~\ref{def:periodic} and $\alpha>0$. Then
$\lim_{R\to\infty}E_{g_\alpha}(C;R)=E^{\mathrm{per}}_{g_\alpha}(C)$. In particular the lower energy $E_{g_\alpha}(C)$ equals
the periodic energy, and the limit exists. \tierC{}
\end{proposition}
\begin{humanproof}
For $x\in C$ let $e(x)=\sum_{y\in C,\,y\ne x}G_\alpha(x-y)$. It is constant on each coset $\Lambda+v_i$, say equal to
$e_i$, and $E^{\mathrm{per}}=\frac1N\sum_ie_i$. Let $n_i(R)$ be the number of points of $\Lambda+v_i$ in
$\overline B(0,R)$. Lattice point counting, done by comparing the disk with the union of fundamental parallelograms
centred at the points it contains, gives $n_i(R)/(\pi R^2)\to1/|\Lambda|$. Hence $n_i(R)/n_C(R)\to1/N$.

\emph{Upper bound.} Dropping the restriction $y\in C_R$ only adds non-negative terms, so
$E(C;R)\le\frac1{n_C(R)}\sum_{x\in C_R}e(x)=\sum_i\frac{n_i(R)}{n_C(R)}e_i\to E^{\mathrm{per}}$.

\emph{Lower bound.} Let $\tau(L)=\sup_{u\in\R^2}\sum_{\lambda\in\Lambda,\,|\lambda+u|>L}G_\alpha(\lambda+u)$. This tends
to $0$ as $L\to\infty$ by Gaussian decay. If $x\in C_{R-L}$, every $y\in C$ with $|x-y|\le L$ lies in $C_R$, so the
inner sum at $x$ is at least $e(x)-N\tau(L)$. Discarding the points of $C_R\setminus C_{R-L}$, which only lowers the sum,
\[
  E(C;R)\ \ge\ \frac{1}{n_C(R)}\sum_{x\in C_{R-L}}\bigl(e(x)-N\tau(L)\bigr)\ \xrightarrow[R\to\infty]{}\ E^{\mathrm{per}}-N\tau(L).
\]
The limit uses $n_i(R-L)/n_C(R)\to1/N$. Letting $L\to\infty$ shows that the $\liminf$ is at least $E^{\mathrm{per}}$.
\end{humanproof}

Together, Theorem~\ref{thm:lp-periodic} and Proposition~\ref{prop:periodic-limit} prove: \emph{for every periodic
density-one configuration $C$ and every $\alpha>0$ for which a sharp auxiliary function exists,
$E_{g_\alpha}(C)\ge\sum_{a\in\Lat\setminus0}e^{-\pi\alpha|a|^2}=E_{g_\alpha}(\Lat)$} \tierC{}.

\subsection{What upstream adds}

Two things are missing from the argument above. Upstream~\cite{oai} supplies both.
\begin{enumerate}
\item \textbf{Existence of sharp auxiliary functions for every $\alpha>0$.} This is the hard part. In dimension two no
construction by modular forms is known; that is the method of CKMRV~\cite{ckmrv} in dimensions 8 and 24. Upstream
instead builds $f_\alpha$ from a sine product with double zeros encoding the periodic contact set. Multiplying by
double-pole and simple-pole terms gives entire functions with prescribed values and slopes. Compactly supported
spectral measures, Gaussian damping, and finite reference lists handled by positive quadrature complete the
construction. The quadrature rests on certified matrix and polynomial inequalities, checked with python-flint ball
arithmetic. \tierL{}~\cite{oai}
\item \textbf{Arbitrary competitors.} The competitor need not be periodic. Upstream compares sums over disk points with
area integrals, with a uniform error, and obtains $E_{g_\alpha}(C)\ge\widehat f_\alpha(0)-f_\alpha(0)$ for every locally
finite configuration of centred density one. \tierL{}~\cite{oai}. The full theorem, including this step, is
kernel-checked in Lean \tierA{} (Theorem~\ref{thm:main-statement}).
\end{enumerate}

% ====================================================================================================
\section{From Gaussians to all completely monotone potentials}
\label{sec:mixtures}

Suppose we know, for every $\alpha>0$ and every admissible competitor $C$,
\begin{equation}\label{eq:gauss-bound}
  E_{g_\alpha}(C)\ \ge\ L(\alpha):=\sum_{a\in\Lat\setminus0}e^{-\pi\alpha|a|^2}.
\end{equation}
How do we get the bound for a general admissible $g$? Here is the human argument in outline \tierC{}.

\emph{Step 1 (Bernstein).} Write $g(t)=\int_{[0,\infty)}e^{-\beta t}\,d\mu(\beta)$ as in Theorem~\ref{thm:bernstein}.
In the normalisation of \eqref{eq:gauss-bound} this reads $e^{-\beta t}=g_{\beta/\pi}(t)$.

\emph{Step 2 (finite sums commute with the integral).} For fixed $R$ the disk energy is a \emph{finite} sum of
non-negative terms, so
\[
  E_g(C;R)=\int_{[0,\infty)}E_{g_{\beta/\pi}}(C;R)\,d\mu(\beta).
\]
\emph{Step 3 (Fatou, not linearity).} The $\liminf$ is not linear: in general
$\liminf(u_R+w_R)\ne\liminf u_R+\liminf w_R$. The correct tool is Fatou's lemma, which gives exactly the inequality we
need. Fatou is a statement about \emph{sequences}, while our $\liminf$ is over a real parameter $R\to\infty$. So first
choose radii $R_n\to\infty$ with $E_g(C;R_n)\to E_g(C)$, and then apply Fatou along $(R_n)$:
\[
  E_g(C)=\lim_n\int E_{g_{\beta/\pi}}(C;R_n)\,d\mu\ \ge\ \int\liminf_n E_{g_{\beta/\pi}}(C;R_n)\,d\mu
  \ \ge\ \int E_{g_{\beta/\pi}}(C)\,d\mu(\beta).
\]
The last step holds because a $\liminf$ along a subsequence of radii is at least the $\liminf$ over all radii.

\emph{Step 4 (the Gaussian bound, then Tonelli).} On $(0,\infty)$ use \eqref{eq:gauss-bound}. On the lattice side,
Tonelli's theorem (non-negative terms) swaps sum and integral:
\[
  \int_{(0,\infty)}L(\beta/\pi)\,d\mu(\beta)=\sum_{a\in\Lat\setminus0}\int_{(0,\infty)}e^{-\beta|a|^2}d\mu(\beta).
\]
\emph{Step 5 (the atom at $0$).} If $\mu(\{0\})=c>0$ then $g\ge c$, so $E_g(C)\ge c\,E_1(C)=\infty$ by the constant-potential
fact of \S\ref{sec:cm}, and the bound is trivial. If $\mu(\{0\})=0$, Steps 3--4 give
$E_g(C)\ge\sum_{a\ne0}g(|a|^2)=\mathcal E_g(\Lat)$.

\emph{Step 6 (attainment).} $\Lat$ is itself locally finite and of density one (kernel-checked in our file
\lean{AnyDensity.lean} as \lean{A_lf} and \lean{A_d1}, which are used in the proof of the audited
\lean{TriangularDensity.attained} \tierA{}), so the bound applies to it. For $C=\Lat$, each point sees at most the full lattice sum, so
$E_g(\Lat;R)\le\mathcal E_g(\Lat)$ for every $R$ (the upper-bound step of Proposition~\ref{prop:periodic-limit}). Then
$\limsup\le\mathcal E_g(\Lat)\le\liminf$, the second inequality being the bound just proved for $C=\Lat$.

\subsection{How upstream completes it}

Upstream completes the passage with lower-limit, endpoint and infinite-energy arguments (upstream manuscript, \S5)
\tierL{}~\cite{oai}. Its Lean proof follows a route that avoids Bernstein's theorem on the open interval and its endpoint
difficulties at $t\to0$. We read the following off upstream's Lean files; the description of the route is ours
\tierC{}. The proof of \lean{universal_energy_minimum} in \lean{Universal.lean} calls
\lean{universal_from_gaussian_bounds}, which is built from the lemmas below. So they lie in the dependency cone of
the audited theorem, whose axiom footprint is the allowlist (\S\ref{sec:lean}). We did not print their axioms
separately.
\begin{itemize}
\item \emph{Shift.} For $\varepsilon>0$ the shifted function $g(\varepsilon+t)$ is a mixture
$\int_{[0,1]}v^t\,d\rho(v)$ of \emph{power potentials} $v^t$. Here $\rho$ is a finite measure on $[0,1]$ of mass
$g(\varepsilon)$, obtained by upstream's own finite-difference argument (Lean: \lean{CM.shifted_mixture}). With
$v=e^{-\pi\alpha}$ one has $v^t=e^{-\pi\alpha t}$: $v\in(0,1)$ are Gaussians, $v=1$ is the constant, and $v=0$ is the
zero potential.
\item \emph{Gaussian to power bounds.} The case split $v=0$ / $v=1$ / $0<v<1$ is the following lemma.
\item \emph{Removing the shift.} Since $g$ is antitone, $E_{g(\varepsilon+\cdot)}(C)\le E_g(C)$
(\lean{shifted_energy_le}). The lattice energy is the supremum of the shifted ones over $\varepsilon=1/(n+1)$
(\lean{latticeEnergy_shift_iSup}). Fatou and Tonelli enter through \lean{shifted_mixture_energy}.
\end{itemize}

\begin{leanbox}[: upstream, \lean{OAI/Analysis/Triangular/Energy/Potentials.lean} (verbatim)]
\begin{lstlisting}[language=Lean]
lemma power_bounds_from_gaussians {C : Set Plane} (hD : DensityOne C)
    (hgauss : ∀ α : ℝ, 0<α → latticeEnergy (gaussPot α)≤energy (gaussPot α) C)
    (v : I) : latticeEnergy (powerPotential v)≤energy (powerPotential v) C := by
  by_cases h0 : (v:ℝ)=0
  · have hl : latticeEnergy (powerPotential v)=0 := by
\end{lstlisting}
(\dots) and, in the case $v=1$:
\begin{lstlisting}[language=Lean]
  by_cases h1 : (v:ℝ)=1
  · have he : powerPotential v=(fun _ => (1:ℝ)) := by
      funext t
      simp only [powerPotential,h1,Real.one_rpow]
    rw [he,energy_one hD]
    exact le_top
\end{lstlisting}
\end{leanbox}

So the atom at $\alpha=0$ ($v=1$) is handled exactly as in Step~5, by \lean{energy_one}.

% ====================================================================================================
\section{A Lean 4 primer for this book}
\label{sec:lean}

You do not need to write Lean to use this book, but you must be able to \emph{read} a statement and know what a
``kernel-checked'' label does and does not certify.

\subsection{Reading a theorem statement}

\begin{leanbox}[: the main theorem as stated in the challenge file (verbatim)]
\begin{lstlisting}[language=Lean]
/-- Universal energy minimality for smooth completely monotone potentials, with attainment. -/
theorem universal_energy_minimum (g : ℝ → ℝ) (C : Set Plane)
    (hg : AdmissiblePotential g) (hC : LocallyFinite C) (hd : DensityOne C) :
    latticeEnergy g ≤ energy g C ∧ latticeEnergy g = energy g A := by
  sorry
\end{lstlisting}
\end{leanbox}

Read it left to right.
\begin{itemize}
\item \lean{(g :} $\R\to\R$\lean{) (C : Set Plane)} are the objects: a potential and a configuration.
\item \lean{(hg : ...) (hC : ...) (hd : ...)} are \emph{hypotheses}. Each is a named proof that the stated property holds.
To apply the theorem you must supply all three.
\item After the final colon comes the \emph{conclusion}. Here it is a conjunction ($\wedge$) of the bound and the
attainment.
\item \lean{:= by sorry} is the proof. In the \emph{challenge} file it is \lean{sorry} \emph{by design}. The proved version
is \lean{OAI/Analysis/Triangular/Energy/Universal.lean}. The Comparator tool certifies that the two statements agree, and
the build of the proved file is kernel-checked.
\end{itemize}

In human language:

\begin{theorem}[Upstream's main theorem~\cite{oai}]\label{thm:main-statement}
Let $g$ be admissible (Definition~\ref{def:cm}) and $C\subseteq\R^2$ locally finite with centred density one. Then
\[
  \mathcal E_g(\Lat)\ \le\ E_g(C)\qquad\text{and}\qquad \mathcal E_g(\Lat)=E_g(\Lat),
\]
with values in $[0,\infty]$, and $E_g$ the lower energy over closed centred disks. \tierA{} for the Lean statement
\lean{OAI.AtomicTriangular.universal_energy_minimum} (kernel-checked, conditional on the fidelity of the definitions).
The proof is AI-generated and not peer reviewed.
\end{theorem}

\subsection{Corollaries are short}

Once the main theorem is available, a corollary only needs to check the hypothesis. Here is the Riesz case, quoted
verbatim from our file \lean{formal_cert/Triangular/Riesz.lean}:

\begin{leanbox}[: \lean{formal_cert/Triangular/Riesz.lean} (verbatim excerpt)]
\begin{lstlisting}[language=Lean]
/-- Riesz potential in the variable `t = r²`. -/
noncomputable def riesz (s : ℝ) : ℝ → ℝ := fun t => t ^ (-(s / 2))

theorem riesz_admissible {s : ℝ} (hs : 0 < s) : AdmissiblePotential (riesz s) := by
  refine ⟨?_, ?_, ?_⟩
  · intro x hx
    exact (Real.contDiffAt_rpow_const_of_ne (ne_of_gt hx)).contDiffWithinAt
  · intro t ht
    exact (Real.rpow_pos_of_pos ht _).le
\end{lstlisting}
\dots
\begin{lstlisting}[language=Lean]
theorem triangular_riesz_optimal {s : ℝ} (hs : 0 < s) (C : Set Plane) (hC : LocallyFinite C) (hd : DensityOne C) :
    latticeEnergy (riesz s) ≤ energy (riesz s) C :=
  (universal_energy_minimum (riesz s) C (riesz_admissible hs) hC hd).1
\end{lstlisting}
\end{leanbox}

The anonymous constructor $\langle$\lean{?_, ?_, ?_}$\rangle$ splits \lean{AdmissiblePotential} into its three
conjuncts: smoothness, non-negativity and complete monotonicity. Each bullet ($\cdot$) proves one of them. The last line applies
the main theorem and takes the first component \lean{.1} of the conjunction. The docstring in the file also remarks that
for $s\le2$ the lattice energy is $\infty$, so the corollary is then vacuous. That remark is a comment, not a theorem
\tierC{}. The \emph{finite} side ($s>2$) is the kernel-checked \lean{latticeEnergy_riesz_lt_top} \tierA{}.

\subsection{\texttt{sorry}, axioms and the allowlist}

Lean accepts a file containing \lean{sorry}. It prints a warning and \emph{exits with code $0$}; the recorded compile
below shows exactly this (return code $0$ with a \lean{sorry} warning). Behind the scenes,
\lean{sorry} is the axiom \lean{sorryAx}, which proves anything. A user can also declare a new \lean{axiom}, for example
one asserting \lean{False}, and prove anything from it; again the compiler exits $0$. The exit code is therefore
\emph{not} a proof check. What is checked is the list of axioms a declaration depends on, printed by \lean{#print axioms}.
This book accepts exactly the allowlist
\[
  \{\texttt{propext},\ \texttt{Classical.choice},\ \texttt{Quot.sound}\},
\]
the standard axioms of Lean's classical mathematics. Anything else, in particular \lean{sorryAx}, rejects the result.

Our audit file \lean{formal_cert/Triangular/Axioms.lean} prints the axioms of 22 declarations and of two planted
controls. The controls, quoted verbatim, are:
\begin{lstlisting}[language=Lean]
theorem ctl_sorry : TriangularRiesz.riesz 1 1 = 0 := by sorry
axiom ctl_cheat : False
theorem ctl_axiom : TriangularRiesz.riesz 1 1 = 0 := ctl_cheat.elim
#print axioms ctl_sorry
#print axioms ctl_axiom
\end{lstlisting}
Both controls ``prove'' the false statement $1^{-1/2}=0$. The recorded compile of this file
(\lean{results/discovery/D1/compile_Axioms.json}, return code $0$) contains, among its output lines:
\begin{lstlisting}
'TriangularRiesz.triangular_riesz_optimal' depends on axioms: [propext, Classical.choice, Quot.sound]
'OAI.AtomicTriangular.universal_energy_minimum' depends on axioms: [propext, Classical.choice, Quot.sound]
.../formal_cert/Triangular/Axioms.lean:32:8: warning: declaration uses `sorry`
'ctl_sorry' depends on axioms: [propext, sorryAx, Classical.choice, Quot.sound]
'ctl_axiom' depends on axioms: [ctl_cheat, propext, Classical.choice, Quot.sound]
\end{lstlisting}
(The path in the warning line is shortened here.) The return code is $0$ even though two false theorems were
``proved''. The audit catches them, because their footprints contain \lean{sorryAx} and \lean{ctl_cheat}. This is the
positive and negative control of the audit: it must flag the planted cheats, and it does.

\subsection{What ``kernel-checked'' attaches to}

Following the Elenchus method~\cite{elenchus}, the label \tierA{} attaches to the \emph{formal statement}, never to its
gloss in English. Lean certifies that \lean{universal_energy_minimum}, \emph{as written with upstream's definitions},
follows from the allowlisted axioms. It does not certify that \lean{DensityOne} is the density a physicist means, or that
\lean{energy} (a $\liminf$ over closed centred disks) is the energy CKMRV mean. Those are questions of \emph{statement
adequacy}, and this chapter's \S\ref{sec:configs} raises several of them: centred disks, closed disks, $\liminf$, and
the empty-set convention of \lean{diskPoints}. No human has audited statement adequacy for these files. The ledger
records 6 open flags. So the correct phrase is ``kernel-checked, conditional on the fidelity of the definitions''.
Avoid ``proved true'', ``certified physics'' or ``100\% verified''.

% ====================================================================================================
\section{Reviewing a draft}
\label{sec:review}

A student draft titled \emph{``Universal Optimality of the Planar Triangular Lattice: A Formal Proof Architecture''}
(header: ``Statement-Adequacy Review Draft'') summarised the material of this chapter. Reviewing it is a good exercise.
Before reading the tables, try to find what is wrong with each quoted claim yourself.

\begin{center}\small
\begin{tabular}{@{}p{0.4cm}p{3.4cm}p{5.4cm}p{4.7cm}@{}}
\toprule
& \textbf{Draft says} & \textbf{Correct statement} & \textbf{Why}\tabularnewline
\midrule
1 & \raggedright ``derivatives strictly alternate in sign'' &
\raggedright $(-1)^rg^{(r)}(t)\ge0$ for all $r\ge0$, $t>0$ (weak inequalities), with $g$ $C^\infty$ on $(0,\infty)$. &
\raggedright Constants are admissible. The Lean definition \lean{AdmissiblePotential} uses $0\le$, not $0<$.
\tabularnewline\addlinespace
2 & \raggedright Bernstein's representation is ``strictly sufficient''; the passage from Gaussians to mixtures is not
justified &
\raggedright $g(t)=\int_{[0,\infty)}e^{-\alpha t}d\mu(\alpha)$ with $\mu\ge0$, and $\mu$ may charge $\alpha=0$. Gaussian
bounds are \emph{sufficient}. The passage needs Tonelli, Fatou along a sequence of radii, and the atom at $0$ treated
separately. &
\raggedright $\liminf$ is not linear (\S\ref{sec:mixtures}). The constant part has infinite energy at density one
(\lean{energy_one}).
\tabularnewline\addlinespace
3 & \raggedright Cohn--Elkies-type lemma by ``majorization'' &
\raggedright $f\le G_\alpha$ (a \emph{minorant}) and $\widehat f\ge0$ give $E(C)\ge\widehat f(0)-f(0)$ at density one.
Root conditions on $\Lat\setminus0$ and $\Lat^*\setminus0$ give equality on $\Lat$ by Poisson summation. Here
$G_\alpha=e^{-\pi\alpha|x|^2}$. &
\raggedright Theorem~\ref{thm:lp-periodic}. A majorant gives no lower bound. The draft also omits the bound itself, the
density, and the normalisation.
\tabularnewline\addlinespace
4 & \raggedright auxiliary functions from ``integral transforms of weakly holomorphic modular forms adapted to hexagonal
symmetry'' &
\raggedright Upstream uses sine-product cardinal interpolation, Gaussian damping and certified finite quadrature. &
\raggedright Modular forms are the CKMRV method in dimensions 8 and 24~\cite{ckmrv}. No modular-form construction is
known in dimension 2.
\tabularnewline
\bottomrule
\end{tabular}
\end{center}

\begin{center}\small
\begin{tabular}{@{}p{0.4cm}p{3.4cm}p{5.4cm}p{4.7cm}@{}}
\toprule
& \textbf{Draft says} & \textbf{Correct statement} & \textbf{Why}\tabularnewline
\midrule
5 & \raggedright ``We demonstrate'' the theorem; author line ``Statement-Adequacy Review Draft'' &
\raggedright The main theorem is upstream's~\cite{oai}: AI-generated, kernel-checked, Comparator-accepted, not peer
reviewed. Our contribution is corollaries and an audit. &
\raggedright Attribution, and honesty about the review status.
\tabularnewline\addlinespace
6 & \raggedright Theorem 2 (arbitrary densities) without the hypotheses of local finiteness and density $\rho$ &
\raggedright For $\rho>0$ and $C$ locally finite with density $\rho$ (\lean{DensityRho}): $E_g(C)\ge\mathcal
E_{g(\rho^{-1}\cdot)}(\Lat)$, attained by $\rho^{-1/2}\Lat$. The energy is the lower energy. &
\raggedright Without local finiteness the counts are meaningless. Lean: \lean{universal_any_density} (D2, see
Chapter~2).
\tabularnewline\addlinespace
7 & \raggedright Theorem 3 (non-uniqueness), with no qualifier on the class of configurations &
\raggedright Non-uniqueness holds among \emph{all} locally finite configurations: $\Lat\setminus\{0\}$ is optimal and is
not an isometric image of $\Lat$. Uniqueness among \emph{periodic} configurations is open in 2D and not claimed. &
\raggedright $\Lat\setminus\{0\}$ is not periodic. Lean:\\{\footnotesize\lean{universal_optimum_not_unique}}.\\ Upstream: ``It identifies
the minimum value, not all minimizers.''
\tabularnewline\addlinespace
8 & \raggedright (missing) &
\raggedright $E_g$ is the \emph{lower} energy ($\liminf$). CKMRV Def.~1.3 also requires the limit to exist. Disks are
closed and centred. &
\raggedright The Lean notion is weaker than CKMRV's. Open vs closed disks: \DFourStatus.
\tabularnewline
\bottomrule
\end{tabular}
\end{center}

The tiers in the tables are those of the sources cited in each row. Rows 1, 6 and 7 point to kernel-checked Lean
declarations \tierA{}. Row 3 rests on Theorem~\ref{thm:lp-periodic} \tierC{} for periodic competitors and on
upstream \tierL{}~\cite{oai} for arbitrary ones. Row 4 is \tierL{}~\cite{oai,ckmrv}. Row 8 is a statement-adequacy
remark \tierC{}. Row 2 combines \tierL{} (Bernstein) and \tierC{} (the Fatou
argument).

\subsection{The corrected statement on one page}

\begin{definition}[Setting]
A configuration $C\subseteq\R^2$ is locally finite. It has density one if $\#(C\cap\overline B(0,R))/(\pi R^2)\to1$
(closed, centred disks). A potential $g$ of the squared distance is admissible if it is $C^\infty$ on $(0,\infty)$ and
$(-1)^rg^{(r)}\ge0$ there for all $r\ge0$. The lower energy is
$E_g(C)=\liminf_{R\to\infty}\frac{1}{\#C_R}\sum_{x\ne y\in C_R}g(|x-y|^2)\in[0,\infty]$.
\end{definition}

\begin{lemma}[LP bound, density one]
Let $\alpha>0$ and $G_\alpha(x)=e^{-\pi\alpha|x|^2}$, $\widehat f(\xi)=\int f(x)e^{-2\pi ix\cdot\xi}dx$. If $f$ is radial
Schwartz with $f\le G_\alpha$ on $\R^2\setminus0$ and $\widehat f\ge0$, then every density-one configuration $C$
satisfies $E_{e^{-\pi\alpha t}}(C)\ge\widehat f(0)-f(0)$. This is \tierC{} for periodic $C$ (Theorem~\ref{thm:lp-periodic}
with Proposition~\ref{prop:periodic-limit}) and \tierL{}~\cite{oai} for arbitrary locally finite $C$. If moreover
$f=G_\alpha$ on $\Lat\setminus0$ and $\widehat f=0$ on $\Lat^*\setminus0$, then
$\widehat f(0)-f(0)=\sum_{a\in\Lat\setminus0}e^{-\pi\alpha|a|^2}=E_{e^{-\pi\alpha t}}(\Lat)$ \tierC{}.
\end{lemma}

\begin{theorem}[Upstream~\cite{oai}; AI-generated, not peer reviewed]
For every admissible $g$ and every locally finite $C$ of density one,
$\sum_{a\in\Lat\setminus0}g(|a|^2)\le E_g(C)$, and equality holds for $C=\Lat$. The values are in $[0,\infty]$. The
theorem identifies the minimum value, not all minimisers. \tierA{} for the Lean statement
\lean{OAI.AtomicTriangular.universal_energy_minimum}, conditional on the fidelity of the definitions. The strong
(per-potential) Cohn--Kumar conjecture is \emph{not} claimed and remains open.
\end{theorem}

% ====================================================================================================
\section{Exercises}
\label{sec:ex1}

Difficulty: (E) easy, (M) medium, (H) hard.

\begin{exercise}[(E)]
Show that $\Z^2$ and $\Lat$ both have density one in the sense of Definition~\ref{def:density}. You may use the lattice
point count $n(R)=\pi R^2/|\Lambda|+O(R)$. Why does Definition~\ref{def:density} need local finiteness?
\end{exercise}

\begin{exercise}[(E)]
Verify the Gram matrix of Proposition~\ref{prop:gram} and find all $(j,k)$ with $j^2+jk+k^2=13$. Check that there are
twelve.
\end{exercise}

\begin{exercise}[(E)]
Which of these are admissible in the sense of Definition~\ref{def:cm}? (a) $e^{-t}+3$; (b) $t^{-1}$; (c) $e^{-t^2}$;
(d) $1/(1+t)$; (e) $\sqrt t$. For (c), compute $g''(t)$ near $t=0$.
\end{exercise}

\begin{exercise}[(E)]
Show that $E_1(C;R)=n_C(R)-1$ for the constant potential $g\equiv1$, and deduce that $E_1(C)=\infty$ at density one.
Find the corresponding upstream Lean lemma and compare its hypotheses with yours.
\end{exercise}

\begin{exercise}[(M)]
Prove that a product of two completely monotone functions is completely monotone. Use this and
\eqref{eq:yukawa-laplace} to give two different proofs that $e^{-\kappa\sqrt t}/\sqrt t$ is completely monotone. (For
the first proof you also need: $e^{-\kappa\sqrt t}$ is completely monotone. Prove it from
\eqref{eq:yukawa-laplace} with an extra differentiation in $\kappa$, or via the composition rule.)
\end{exercise}

\begin{exercise}[(M)]
Prove Proposition~\ref{prop:gauss-ft} in one variable by showing that $\phi(\xi)=\int e^{-\pi\alpha x^2}e^{-2\pi ix\xi}dx$
satisfies $\phi'(\xi)=-(2\pi\xi/\alpha)\phi(\xi)$.
\end{exercise}

\begin{exercise}[(M)]
Repeat the proof of Theorem~\ref{thm:lp-periodic} at density $\rho=N/|\Lambda|$ and show that the bound becomes
$\rho\widehat f(0)-f(0)$. Where exactly does the density enter?
\end{exercise}

\begin{exercise}[(M)]
Give an example of real sequences with $\liminf(u_n+w_n)>\liminf u_n+\liminf w_n$. Explain why Step~3 of
\S\ref{sec:mixtures} would be wrong if it claimed equality, and why the inequality that Fatou gives is the one needed.
\end{exercise}

\begin{exercise}[(M)]
Let $C=\Lat\setminus\{0\}$. Show directly that $C$ has density one, and that $E_g(C)=E_g(\Lat)$ for $g$ Gaussian. Why is
$C$ not periodic? Which Lean declaration states the general version?
\end{exercise}

\begin{exercise}[(H)]
Complete the lattice point counting step of Proposition~\ref{prop:periodic-limit}: show that
$n_i(R)=\pi R^2/|\Lambda|+O(R)$ for a translate $\Lambda+v_i$, uniformly in $v_i$.
\end{exercise}

\begin{exercise}[(H)]
Suppose a radial Schwartz $f$ satisfies (i)--(iv) of Definition~\ref{def:aux}. Show that $\nabla(G_\alpha-f)=0$ at every
$a\in\Lat\setminus0$. In one radial variable, explain why the roots of $G_\alpha-f$ at $|a|$ have even order. Why does
this suggest a construction from a product with double zeros?
\end{exercise}

\begin{exercise}[(H)]
The draft review (\S\ref{sec:review}) lists eight problems. Find a \emph{ninth} adequacy question about upstream's
definitions that the table does not cover. Hint: think about \lean{diskPoints} on a configuration that is not locally
finite, or about centred versus uncentred disks. State it precisely and say which tier an answer could reach.
\end{exercise}

\begin{exercise}[(M, Python)]\label{ex:py-theta}
\emph{Poisson summation and the theta comparison.} The script below computes truncated theta sums
$\theta_\Lat(\alpha)$, the Poisson-dual expression $\theta_\Lat(1/\alpha)/\alpha$ (Example~\ref{ex:theta}), and
$\theta_{\Z^2}(\alpha)$ for the square lattice, which also has density one. It was run and printed the output below
\tierX{}.
\begin{lstlisting}[language=Python]
import math

B = math.sqrt(3) / 2  # b = sqrt(3)/2, as in upstream


def norm_sq(j: int, k: int) -> float:
    """|p(j,k)|^2 = (1/b)(j^2 + jk + k^2) for the covolume-one triangular lattice A."""
    return (j * j + j * k + k * k) / B


def theta_A(alpha: float, M: int = 40) -> float:
    """Truncated theta sum over A: sum of exp(-pi*alpha*|a|^2), |j|,|k| <= M."""
    return sum(math.exp(-math.pi * alpha * norm_sq(j, k))
               for j in range(-M, M + 1) for k in range(-M, M + 1))


def theta_Z2(alpha: float, M: int = 40) -> float:
    """Theta sum over the square lattice Z^2 (also covolume one)."""
    s = sum(math.exp(-math.pi * alpha * n * n) for n in range(-M, M + 1))
    return s * s


print("alpha   theta_A(alpha)      theta_A(1/alpha)/alpha   theta_Z2(alpha)")
for alpha in (0.5, 1.0, 2.0):
    lhs = theta_A(alpha)
    rhs = theta_A(1.0 / alpha) / alpha
    print(f"{alpha:4.1f}  {lhs:.15f}  {rhs:.15f}  {theta_Z2(alpha):.15f}")
\end{lstlisting}
\begin{lstlisting}
alpha   theta_A(alpha)      theta_A(1/alpha)/alpha   theta_Z2(alpha)
 0.5  2.008477918597919  2.008477918597920  2.014967440690170
 1.0  1.159595266963928  1.159595266963928  1.180340599016097
 2.0  1.004238959298960  1.004238959298960  1.007483720345085
\end{lstlisting}
(a) Explain why columns 2 and 3 agree, citing the two propositions and the theorem used. Why do they agree only to
about $10^{-15}$? (b) Compute $\theta_\Lat(1)-1$ and compare with Example~\ref{ex:shells}. (c) The Gaussian energies at
$\alpha=1$ are $\theta-1$: about $0.15960$ for $\Lat$ and $0.18034$ for $\Z^2$ \tierX{}. This is consistent with the main
theorem, but it supports nothing (Tier X). Which theorem in the literature compares theta functions of all lattices of
covolume one? (See~\cite{montgomery}.) (d) Change \texttt{M} to $5$ and explain why the output barely changes.
\end{exercise}

\begin{exercise}[(E, Python)]
Modify the shell-count script of Example~\ref{ex:shells} to list the values $n\le 50$ represented by $j^2+jk+k^2$, and
the corresponding $r(n)$. Conjecture which primes $p$ occur, then check your conjecture by hand for $p\le 13$. Report
only what your run prints.
\end{exercise}

\begin{exercise}[(H)]
Using Theorem~\ref{thm:lp-periodic}, show that among \emph{lattices} of covolume one, the triangular lattice minimises
$\sum_{\lambda\ne0}e^{-\pi\alpha|\lambda|^2}$ for every $\alpha$ for which a sharp auxiliary function exists. Explain why
this lattice statement is much weaker than Theorem~\ref{thm:main-statement}.
\end{exercise}
